% 영문 경력기술서 조판 선언 (통합 입사지원서의 뒷부분).
%
% pandoc 기본 템플릿이 fontspec 으로 라틴 활자를 세운 뒤 여기가 끼어든다. xetexko 가
% 스크립트를 갈라 한글에만 한글 조판 규칙(줄바꿈·자간)을 적용한다. 이게 없으면 한글이
% 라틴 규칙으로 짜여 어절 중간이 아니라 공백에서만 끊기고, 오른쪽 끝이 심하게 들쭉날쭉해진다.
\usepackage{xetexko}
\setmainhangulfont{Pretendard}[Ligatures=TeX]
\setsanshangulfont{Pretendard}

% 이 문서는 영문이라 xetexko 가 거의 개입하지 않는다. 다만 스마트 따옴표(U+2019)만은
% CJK 문장부호로 잡혀 «runtime' s» 처럼 벌어지므로, build.sh 에서 `-f markdown-smart` 로
% 끄고 ASCII 어포스트로피를 그대로 쓴다. 표제와 이름에 섞인 한글은 위 설정이 받는다.

% 한자는 Pretendard 에 글리프가 없다. 본문에 한자를 쓰지 않지만, 인용문에 섞여 들어와도
% PDF 에 빈칸으로 찍히지 않도록 대체를 잡아 둔다.
\setmainhanjafont{D2Coding}

\usepackage[dvipsnames]{xcolor}
\definecolor{accent}{HTML}{1B3A57}
\definecolor{muted}{HTML}{5A6672}

% 표제 위계. 절이 일곱 개라 번호를 크게 두지 않고 색으로만 가른다.
\usepackage{titlesec}
\titleformat{\section}{\normalfont\large\bfseries\color{accent}}{}{0pt}{}
\titlespacing*{\section}{0pt}{12pt plus 2pt minus 2pt}{5pt}

% 들여쓰기 대신 문단 사이 여백. 첫 줄 들여쓰기와 여백을 같이 쓰면 문단 경계가 두 번 표시된다.
%
% 행간은 국문판(1.14)보다 좁다. 라틴 본문은 그 행간이면 성글어 보이고, 무엇보다 1.14 로는
% 마지막 쪽에 불릿 셋만 남아 그 쪽의 9할이 빈다 — 읽는 사람에게는 «5쪽짜리» 로 보이고
% 실제로는 4쪽 반이다. 문장을 깎아 맞추는 대신 판면에서 먼저 찾는다.
\setlength{\parindent}{0pt}
\setlength{\parskip}{3.6pt plus 1.2pt minus 0.4pt}
\linespread{1.06}

% 판면을 몇 pt 넘긴 줄 하나 때문에 문장을 고쳐 쓰는 일이 없게 마지막 줄에 신축을 허용한다.
\setlength{\emergencystretch}{1.2em}

% 머리말 블록(연락처)을 본문 활자보다 작게, 색을 죽여 둔다.
\usepackage{quoting}
\quotingsetup{vskip=6pt,leftmargin=0pt,rightmargin=0pt,indentfirst=false}
\renewenvironment{quote}{\begin{quoting}\small\color{muted}}{\end{quoting}}
